% ============================================================================
% LIMITATIONS AND SCOPE SECTION (Paper-Ready)
% ============================================================================
% Status: FINAL - Reviewer-fest, ehrlich, schützend ohne zu schwächen
% Date: 2025-12-28
% ============================================================================

\section{Limitations and Scope}

\subsection{Scope of the Claims}

The present work is explicitly \textbf{theory-driven} rather than benchmark-oriented. Our goal is not to provide a statistically representative estimate of knowledge activation performance across all language models, domains, or tasks. Instead, we aim to isolate and demonstrate the existence of a specific functional limitation---the gap between latent knowledge and its activation during generative decision-making---under maximally controlled conditions.

Accordingly, our claims are \textbf{existential and diagnostic}, not population-level generalizations. The Activation Challenge is designed to show that such a gap \emph{can} persist, even when ambiguity is minimized and epistemic rules are made explicit.


\subsection{Limited Item Set}

We deliberately use a small, carefully constructed set of ten psychological phenomena. These items are not intended as a random or representative sample of psychological knowledge. Rather, they function as \textbf{diagnostic probes}: linguistically plausible labels that are epistemically demanding due to their proximity to real concepts without being established terminological units.

While a larger item set could increase robustness, it would not change the logical force of our core finding. A single systematic failure under controlled conditions suffices to demonstrate non-equivalence between knowledge possession and knowledge activation.


\subsection{Limited Model Coverage}

Our evaluation includes a small number of contemporary large language models. This choice reflects a trade-off between breadth and interpretability. We do not claim that the observed activation limitations apply uniformly to all models or architectures. Instead, we show that models with comparable latent knowledge can exhibit markedly different activation competence, indicating that $\alpha$ is not reducible to model size, training data volume, or general instruction-following ability.

Future work may expand the model set to explore architectural or training-related correlates of activation competence, but such analysis lies beyond the scope of the present study.


\subsection{Dependence on Epistemic Ground Truth}

Our evaluation relies on expert judgment regarding whether a given label constitutes an established \emph{terminus technicus} in the scientific literature. While this introduces an element of human adjudication, the task is deliberately framed to minimize subjectivity by focusing on \textbf{label existence} and \textbf{terminological establishment}, rather than on theoretical interpretation or empirical effect size.

Importantly, potential disagreements about borderline cases would, if anything, bias results \emph{against} our hypothesis, making false-positive REAL classifications harder---not easier---to detect.


\subsection{Prompt-Specificity and Generalization}

The ESL prompting framework is intentionally restrictive. It operationalizes explicit epistemic stop rules rather than naturalistic conversational norms. As such, our findings should not be interpreted as claims about everyday prompt usage or typical human--AI interaction scenarios.

At the same time, the failure of certain models to improve under these restrictive conditions is precisely the point of interest. The results demonstrate that even maximally explicit epistemic constraints may fail to induce reliable knowledge activation, highlighting a limitation that cannot be dismissed as prompt ambiguity or user error.


\subsection{Absence of Mechanistic Explanation}

Our analysis is functional and behavioral. We do not attempt to identify internal representations, circuits, or neuron-level mechanisms responsible for activation failures. While recent work (e.g., mechanistic interpretability and neuron-level analyses) may offer complementary insights, our contribution lies in establishing the phenomenon at the behavioral level and providing a formal framework in which mechanistic explanations can later be situated.


\subsection{Summary of Limitations}

In summary, this work is limited in scope by design. It does not aim to be exhaustive, statistically representative, or mechanistically complete. Instead, it provides a \textbf{minimal, controlled demonstration} of a previously under-theorized failure mode in large language models. We view this as a necessary first step toward a more comprehensive theory of knowledge activation in generative systems.


% ============================================================================
% INTERNAL NOTES (nicht für Paper)
% ============================================================================
%
% Diese Sektion erreicht:
% - Nimmt Reviewer-Wind aus den Segeln
% - Positioniert als theoretische Isolationsarbeit
% - Erklärt warum klein ≠ schwach
% - Lädt explizit zu Anschlussarbeiten ein
%
